\section{Conclusion}

We study what governs ripple effects in LLM knowledge updating.
Using \textsc{RippleEval}, a verified factual graph from real-world
Wikipedia, we trace how a single fine-tuning update affects connected
facts across hop distances. Errors propagate far beyond the edited
fact, with measurable effects even past five hops in four open-weight
LLMs.

Counterintuitively, \emph{popularity}—the in-degree of an entity in
the factual graph—is a stronger driver of ripple effects than
semantic similarity or long-tail fragility. Popular knowledge is
more easily overwritten as a direct target, more often corrupted as
a neighbor, and more effective at propagating errors to distant
facts. An attention probe further shows that popular knowledge
participates in stronger attention interactions, suggesting a
pathway for this propagation.

Building on this finding, we propose a popularity-aware preservation
strategy that anchors a small set of popular facts via a KL
regularizer. It substantially reduces long-range ripple effects and
consistently outperforms popularity-agnostic baselines, showing that
the property driving instability can also stabilize updating.
Together, our results identify structural connectivity as a unifying
lens on ripple effects, and the released graph may support future
work on knowledge updating and distortion in LLMs.